\subsection{Completed controls for learned weights}

The learned-versus-random comparison conditions on the same coordinates, sequence and graph-statistic features. Negative and inconclusive contrasts remain alongside the positive strict INS/DEL result (\cref{tab:s-hr}). The original $C+S+T$ versus $C+S$ contrasts answer a different question and cannot substitute for these controls.

\begin{table}[H]
\centering
\caption{\textbf{Original-v1 trained-minus-random controls after $C+S+H$.} Each row uses 15 paired fold/seed runs. Exact sign flips use five seed-averaged fold differences; $q$ is the supplied within-metric BH adjustment. These are exploratory comparisons on previously inspected folds.}
\label{tab:s-hr}
\footnotesize
\begin{tabular}{@{}llrlrr@{}}
\toprule
Task & Context & AP difference & 95\% CI & $p$ & $q$ \\
\midrule
\csname @@input\endcsname hr_completed_rows.tex
\bottomrule
\end{tabular}
\end{table}

\subsection{Matched three-class SV sensitivity}

The matched cohort contains 223 events of each type, restricted to equal chromosome/length-bin support. Its macro $C+S$ AP is 0.414145 and $C+S+H$ AP is 0.473760. The latter exceeds $C+S+T$ (0.437789 strict; 0.440860 one-hop), and adding $T$ after $H$ is inconclusive (\cref{tab:s-sv-matched}). This population excludes 75 primary-chromosome inversions lacking matched support; all 298 inversions remain included in the natural-frequency experiment. Absolute AP across these different prevalences and feature constructions should not be used to rank model versions.

\begin{table}[H]
\centering
\caption{\textbf{Three-class common-support macro AP increments.} All 30 runs and 1,170 feature/class evaluations completed. Pointwise bootstrap intervals accompany exact fold sign-flip tests and the supplied within-metric BH adjustment.}
\label{tab:s-sv-matched}
\footnotesize
\begin{tabular}{@{}llrlrr@{}}
\toprule
Conditioning features & Context & $\Delta_T$ & 95\% CI & $p$ & $q$ \\
\midrule
\csname @@input\endcsname sv_matched_rows.tex
\bottomrule
\end{tabular}
\end{table}

\subsection{Additional transfer tasks retain negative and inconclusive findings}

The prospective HG008 refit experiment completed all 30 planned evaluations on 69 clonal INS/DEL variants without training on HG008. Its gains beyond $C+S$ were 0.032525 (\CI $-0.053371$ to 0.110254) strict and 0.064863 ($-0.066739$ to 0.212115) one-hop. Both are inconclusive; a single external genome does not establish generalization across people. Six failed attempts to reproduce historical probes remain distinct from these successful prospective refits.

TraitGym locus-prior prediction \citepS{supp__benegas2025traitgym} retained all 14,780 matched variants in the five-fold adaptation. Neither the original frozen-feature study nor the published NT-2.5B allele-score sensitivity established a positive topology interval. A subsequent completed 60-run probe study improved Mendelian $C+S+V$ AP from 0.102110 with the fixed linear classifier to 0.171056 with fixed shallow boosting, a paired difference of 0.068946 (\CI 0.017423--0.145067), where $V$ denotes published allele scores. This is a classifier improvement on previously inspected data; every topology AP interval remained nonpositive or inconclusive. The shared segment representation is a locus prior, not an allele-specific embedding, and its end-sampled sequence input omits most tested variant bases.

Prediction of measured COSIGT locus quality \citepS{supp__bolognini2026cosigt} retained all 265 loci across 30 runs and 900 evaluations. Every original topology mean-absolute-error interval spanned zero. On the primary mean quality-value fraction, the constant training-median reference had lower mean error than every fitted feature model. A validation-only fallback to this constant reduced observed error but did not establish a topology benefit. These data concern quality conditional on released, evaluable sample--locus pairs; they do not demonstrate improved genotype calling or general callability. Detailed contrasts and the original negative results are retained in \hyperref[sec:s-transfer-boundaries]{Supplementary Information, external and measured-outcome transfer boundaries}.


\subsection{External and measured-outcome transfer boundaries}
\label{sec:s-transfer-boundaries}

\begin{table}[H]
\centering
\caption{\textbf{Retained external-genome, variant-prior and measured-quality results.} Each contrast adds original frozen $T$ to $C+S$ in 15 paired fold/seed runs per context. HG008 uses prospective refits; TraitGym uses the original fixed linear locus-prior adaptation; COSIGT uses the primary continuous QV-fraction target. Positive values favor adding $T$, but AP gains and MAE reductions have different units and cannot be compared in magnitude.}
\label{tab:s-transfer-boundaries}
\footnotesize
\begin{tabular}{@{}lllr l@{}}
\toprule
Endpoint & Context & Metric & Difference & 95\% CI \\
\midrule
\csname @@input\endcsname transfer_boundary_rows.tex
\bottomrule
\end{tabular}
\end{table}

The original TraitGym NT input retained only 22.95\% of tested complex-trait bases and 27.75\% of Mendelian bases before tokenization. This raw-sampling audit does not prove that every retained base survives tokenization or identify the cause of poor prediction. In the allele-score sensitivity, the published NT-2.5B score alone achieved Mendelian AP 0.160370, while $C+S+V$ achieved 0.102110 versus 0.120840 for $C+S$. Its paired $V$ increment was $-0.018730$ (\CI $-0.028974$ to $-0.008670$); the unsuccessful concatenation is retained alongside the later classifier sensitivity. No topology AP contrast in that sensitivity had a wholly positive interval.

For COSIGT, the training-median primary-target MAE was 0.046169, compared with $C+S$ 0.053351 and $C+S+T$ 0.052699 strict/0.053824 one-hop. The later validation-only fallback yielded $C+S+T$ MAE 0.047911/0.047193 and retained the original constant reference. The paired reductions from this fallback were inconclusive in both contexts. Predicting published locus quality is distinct from altering COSIGT calls, measuring per-variant genotype concordance or establishing an all-callable denominator.

\subsection{Masked sequence-feature learning passes development controls; chromosome replication remains pending}
\label{sec:masked-development}

To test a different self-supervised target, we adapted GraphMAE-style masked-feature reconstruction \citepS{supp__hou2022graphmae} to the existing graph encoder. We denote this 48-dimensional, sequence-conditioned representation by $E$ to distinguish it from the original topology-pretrained $T$. Twelve pretraining runs and 24 frozen biological probes (96 feature evaluations) compared full and coordinate-only encoders with matched frozen random controls across three seeds. All development predictions replayed and all probes converged. On fold-A validation chromosomes, mean trained-minus-random AP after $C+S+H$ was 0.005609 for INS/DEL and 0.001836 for cCRE, positive at every seed. Full-trained minus coordinate-only-trained differences were also positive at every seed, averaging 0.006145 and 0.002372 (\cref{tab:s-development}). The coordinate-only ablation also removes model capacity, so this is not an isolated causal estimate of message passing.

These three initializations share one development chromosome block and receive no chromosome confidence interval. The earlier junction candidate failed its full three-seed development gate and remains a failed comparison. The fixed 120-probe chromosome replication and separately prescribed 30-probe original-v1 reference comparison have no completed aggregate result at this evidence cutoff; no partial test metrics enter this manuscript. Fold B must be reported separately because its test chromosomes supplied fold-A development validation. The remaining folds also retain historical v1 label exposure. Development success therefore does not establish superiority to v1, independent external transfer, or a validated general-purpose foundation model.

\subsubsection{All-seed masked-feature development controls}

\begin{table}[H]
\centering
\caption{\textbf{Sequence-conditioned $E$ development controls.} AP values and differences condition on $C+S+H$, using fold-A validation only. All four encoder arms and all three seeds were retained; the table gives the two contrasts discussed in Results. Coordinate-only models have less capacity. Seeds are initializations on the same genomic block, so no chromosome confidence interval is reported.}
\label{tab:s-development}
\footnotesize
\setlength{\tabcolsep}{4pt}
\begin{tabular}{@{}llrrrr@{}}
\toprule
Seed & Task & Full trained & Full random & Trained--random & Full--coordinate \\
\midrule
\csname @@input\endcsname development_rows.tex
\bottomrule
\end{tabular}
\end{table}

The same-width junction-$Q$ reference comparison is also retained: mean $E-Q$ AP after $C+S+H$ was 0.001877 INS/DEL and 0.000099 cCRE, whereas after $C+S$ it was 0.006154 INS/DEL and $-0.000145$ cCRE. The masked-feature and junction experiments used different epoch budgets (20 versus 10); these differences do not isolate the pretraining objective or establish universal improvement. F1 and balanced accuracy were not uniformly better at every seed. The frozen original-v1 comparison and full chromosome replication remain separate pending analyses.

All numerical table rows in this update are generated from bundled completed-study CSV snapshots with recorded source paths and SHA-256 identities. The evidence bundle preserves the complete AP contrast families used here, original intervals and null results. It contains no partially completed chromosome-test scores.
